

\section{Related Work}


\begin{figure*}[!t]
\centering
\includegraphics[width=1\textwidth]{images/framework.pdf} \\
\vspace{-3mm}
\caption{\textbf{Overview of our method.}
There are two-stage training in our framework. In stage I, we leverage a pretrained VGGT model to get a dense scene layout point cloud. Next, the geometry of the 3D Gaussians is supervised by normal and semantic priors distilled from a pretrained model. In state II, we perform inverse rendering based on the Gaussians obtained from the first stage. A mixture light model together with a deshadowing module are employed to model illumination and occlusion relationships, which are then integrated with the Gaussian-rendered G-buffer for physically based (PBR) novel view rendering. Moreover, to ensure consistent material representations during training, we incorporate an additional material consistency constraint.}
\label{fig:pipeline}
% \vspace{-0.4cm}
\end{figure*}

% \subsection{Geometry Reconstruction from Sparse Views}
\textbf{Geometry reconstruction from sparse views.} 
 Existing sparse-view geometry reconstruction methods can be broadly categorized into two groups: per-scene optimization methods and pre-training based methods. Per-scene optimization methods ~\cite{li2024dngaussian,zhu2024fsgs,wang2024use,pooleDreamFusionTextto3DUsing2022,chung2024depth,zhang2024cor,paul2024sp2360,peng2024structure,han2024binocular} focus on learning to recover geometry of a specific scene, and thus cannot generalize to other scenes. In contrast, by taking full advantage of powerful priors from pre-trained models, pretraining based methods allow to generalize beyond the training scenes \cite{wang2024dust3r,wang2025vggt,leroy2024grounding}. Despite the effectiveness of pretraining based methods, they basically lack the ability to achieve material-illumination disentanglement. 
 

% \subsection{Intrinsic Image Decomposition and Inverse Rendering}
\vspace{0.5em}
\noindent \textbf{Inverse Rendering.}
Inverse rendering has long been a fundamental problem in computer vision and computer graphics.
Early methods focus on predicting intrinsic properties from a single image. Representative works include ~\cite{10.1145/3641519.3657445,careaga2024colorful,liang2025diffusion,chen2024intrinsicanything,zeng2024dilightnet,careaga2023intrinsic,kocsis2024intrinsic}. However, intrinsic decomposition from a single image inherently lacks multi-view consistency and often struggles with out-of-distribution cases \cite{du2024gs}.
Hence, recent methods \cite{du2024gs,liang2024gs,gu2025irgs,shi2025gir,wynn2023diffusionerf,ni2025decompositional} aim to recover a 3D-consistent intrinsic field from multiple input images. Although these methods demonstrate impressive results, they focus primarily on object-level scenes, and fail to work well under sparse view conditions involving complex indirect lighting and occlusion. 




